\section*{Chapter 6: Black Holes: Champions Of Curvature}
\begin{description}[style=multiline,leftmargin=2em,labelwidth=1.5em]

% DOCX paragraph 2057
\item[1.]\hypertarget{note:06-01}{}At least for some time, the name \emph{``black hole''} also had a rival in
``\emph{frozen star''}. The latter was coined to emphasise that when a star collapses
and gives life to a black hole, it seems to ``freeze'' in time for an external observer.
However, the rivalry between the two names did not last long, and the scientific
community soon adopted the form we all know today, which is as effective as it is
evocative.

% DOCX paragraph 2058
\item[2.]\hypertarget{note:06-02}{}Conversely, if the black hole had a mass of the order of a few solar masses, the tidal
forces\index{tidal forces} (i.e., the gravitational forces generated by the rapid change
of the gravitational field between one point and another) would be enormously different
even over the length-scale between your head and your feet. Hence, you would be
stretched like a rubber band, ending up like ``spaghetti''. This is an experience that I
frankly do not recommend, not even as a gedanken experiment~\ldots{}

% DOCX paragraph 2059
\item[3.]\hypertarget{note:06-03}{}J. Michell, in \emph{Philosophical Transactions of the Royal Society of London}, 74, 35
(1784).

% DOCX paragraph 2060
\item[4.]\hypertarget{note:06-04}{}R. P. Kerr, in \emph{Physical Review Letters}, 11, 237 (1963).

% DOCX paragraph 2061
\item[5.]\hypertarget{note:06-05}{}R. Ruffini and J. A. Wheeler, in \emph{Physics Today} 24, 30 (1971).

% DOCX paragraph 2062
\item[6.]\hypertarget{note:06-06}{}S. W. Hawking, in \emph{Communications of Mathematical Physics}, 43, 199 (1975); S. W.
Hawking, in \emph{Nature} 248, 3031 (1974).

% DOCX paragraph 2063
\item[7.]\hypertarget{note:06-07}{}As mentioned, supermassive black holes have gravitational fields that vary very
gradually compared to the human scale, or, equivalently, when compared with ``objects''
having dimensions of the order of a couple of metres. That's why I chose one for my
example.

% DOCX paragraph 2064
\item[8.]\hypertarget{note:06-08}{}The two characteristics are not necessarily linked: there are black-hole solutions that
have an event horizon but do not have a physical singularity\index{singularity!physical} at their
centre. For this reason, they are called \emph{regular black holes}.

% DOCX paragraph 2065
\item[9.]\hypertarget{note:06-09}{}P. O. Mazur and E. Mottola, in \emph{Proceedings of the National Academy of Sciences,
USA}, 101, 9545 (2004).

% DOCX paragraph 2066
\item[10.]\hypertarget{note:06-10}{}C. Misner and J. A. Wheeler, in \emph{Annals of Physics}, 2, 525 (1957).

\end{description}
